\documentclass[11pt]{article}

\usepackage[utf8]{inputenc}
\usepackage[T1]{fontenc}
\usepackage{lmodern}
\usepackage{geometry}
\usepackage{amsmath,amssymb,amsfonts,amsthm,mathtools}
\usepackage{bm}
\usepackage{enumitem}
\usepackage{microtype}
\usepackage{hyperref}
\usepackage{titlesec}
\usepackage{booktabs}
\usepackage{array}
\usepackage{setspace}
\usepackage{mathrsfs}
\usepackage{physics}

\geometry{margin=1in}
\hypersetup{colorlinks=true, linkcolor=black, urlcolor=blue, citecolor=black}
\setlist[enumerate]{topsep=0.4em,itemsep=0.35em}
\setlist[itemize]{topsep=0.3em,itemsep=0.25em}
\titleformat{\section}{\large\bfseries}{\thesection.}{0.5em}{}
\titleformat{\subsection}{\normalsize\bfseries}{\thesubsection.}{0.5em}{}
\setstretch{1.06}

\newcommand{\Aether}{\AE{}ther}
\newcommand{\Aflow}{\AE{}ther-flow}
\newcommand{\TheoryName}{\Aflow{} Interpretation of Relativity}
\newcommand{\TheoryProgram}{\Aether{} / \Aflow{} framework}
\newcommand{\Order}{\mathcal{S}}
\newcommand{\Sub}{\mathcal{A}}
\newcommand{\BridgeMetric}{\mathfrak{g}}

\newtheorem{definition}{Definition}
\newtheorem{proposition}{Proposition}
\newtheorem{remark}{Remark}

\title{\textbf{The \TheoryName{}}\\[0.4em]
\large Substrate Kinematics and the Derivational Bridge}
\author{}
\date{}

\begin{document}

\maketitle

\begin{abstract}
This manuscript states the next substantive step beyond the current exact-closure sequence of \TheoryName{}. Its purpose is not to claim that Einsteinian gravity has already been derived from substrate microphysics. Its purpose is narrower and more disciplined: to make explicit the minimum substrate kinematics needed for a future derivational program.

The paper therefore introduces a four-dimensional substrate manifold \(\Sub\) equipped with a positive-definite substrate metric \(G_{AB}\), a unit ordered-flow field \(U^A\), an order scalar \(S(X)\), and additional substrate variables \(Q\). It then defines the local experiential three-space by the projector orthogonal to \(U^A\), constructs a Lorentzian bridge metric \(\BridgeMetric_{AB} = G_{AB} - 2 U_A U_B\), and uses a local observer map to recover the effective relativistic metric \(g_{\mu\nu}\) seen by embedded observers.

The resulting claim is kinematic rather than fully dynamical. The paper states the missing bridge from ontology to explicit substrate structure, shows how observed three-dimensional space and S-time can be anchored in that structure, and records the conditions any completed substrate action must satisfy if it is to recover the exact Einsteinian closure already adopted elsewhere in the active sequence. Exact closure therefore remains the benchmark and reversion theory until those stronger derivational conditions are actually met.
\end{abstract}

\section{Introduction}

The active manuscript sequence of \TheoryName{} now fixes the ontology, the exact effective dynamics, the consistency conditions, the relativistic recovery statement, and the congruence-based flow geometry of the exact-closure theory. What it still lacks is the deeper bridge from ontology to explicit substrate kinematics and derivation.

That missing step is the subject of the present paper. The goal is not another GR / SR recovery statement. Nor is the goal to replace the exact Einsteinian closure already adopted in the active theory line. The goal is to define the minimum substrate-level data from which a future derivation would even have to begin.

\paragraph{Framework and claim boundary.}
\TheoryProgram{} denotes the broader \Aether{} / \Aflow{} framework built on the claims that the \Aether{} is the underlying four-dimensional substrate of reality, the \Aflow{} is its intrinsic ordered motion, observed three-dimensional space is the local experiential slice of that deeper substrate, S-time is the experienced order of change arising from matter, light, and the \Aflow{}, and observed expansion is the three-dimensional appearance of deeper four-dimensional ordered motion. Within that framework, \TheoryName{} denotes the exact-closure theory in which the effective gravitational dynamics are adopted to be exactly Einsteinian with universal matter coupling. In the active sequence, \emph{adoption} means use of that established relativistic dynamics without claiming substrate derivation, while \emph{derivation} is reserved for a first-principles recovery from explicit substrate variables.

The present manuscript stays strictly within that claim boundary. It does not claim a completed recovery of the Einstein equations from substrate variables. It does not present a finished microphysical action that has already passed a consistency analysis. It does make one positive move: it replaces the previously schematic phrase ``future substrate manuscript'' with an explicit kinematic bridge whose success conditions can be stated cleanly.

\section{Role of the Derivational-Bridge Step}

The foundations manuscript already introduced a four-dimensional substrate manifold \(\Sub\), an ordered-motion map \(\Phi_\lambda\), and the idea that observed three-dimensional space is a local experiential slice of the deeper substrate. The dynamics manuscript then fixed the exact effective law of \TheoryName{} to be Einsteinian. The flow-geometry manuscript identified the \Aflow{} with a congruence-based interpretation of that exact relativistic geometry.

What is still missing between those layers is the explicit substrate side of the bridge. In particular, the theory still requires:
\begin{enumerate}
    \item a disciplined specification of the substrate variables from which a relativistic metric sector could emerge;
    \item an explicit local bridge from those substrate variables to the effective Lorentzian geometry seen by observers;
    \item a clear statement of what counts as a successful derivation and what would still count only as an unfinished ansatz.
\end{enumerate}

The present manuscript addresses those three points. It should therefore be read as a kinematic bridge paper, not as a completed derivation paper.

\section{Minimal Substrate Kinematics}

\begin{definition}[Minimal substrate data]
The minimal kinematic data of the derivational bridge are the tuple
\[
\bigl(\Sub, G_{AB}, U^A, S, Q\bigr),
\]
where:
\begin{enumerate}
    \item \(\Sub\) is a smooth four-dimensional substrate manifold;
    \item \(G_{AB}\) is a positive-definite substrate metric on \(\Sub\);
    \item \(U^A\) is a smooth unit ordered-flow field satisfying
    \begin{equation}
    G_{AB} U^A U^B = 1;
    \label{eq:unitflow}
    \end{equation}
    \item \(S:\Sub \to \mathbb{R}\) is an order scalar satisfying
    \begin{equation}
    U^A \partial_A S > 0;
    \label{eq:ordermonotone}
    \end{equation}
    \item \(Q\) denotes any additional substrate variables needed for a later completed dynamics.
\end{enumerate}
\end{definition}

Equation \eqref{eq:unitflow} makes \(U^A\) the explicit kinematic carrier of \Aflow{} at the substrate level. Equation \eqref{eq:ordermonotone} makes \(S\) an order variable that increases along the ordered motion.

This separation between \(U^A\) and \(S\) is important. The ordered-flow field and the order scalar should not be collapsed automatically into one object. If they were identified too early, the framework could accidentally forbid the rotational and vorticity sectors already needed in the relativistic flow-geometry manuscript.

The ordered-motion map introduced in the foundations manuscript is now made explicit through the integral curves of \(U^A\):
\begin{equation}
\dv{X^A}{\lambda} = U^A\bigl(X(\lambda)\bigr),
\qquad
\Phi_\lambda(X_0) = X(\lambda).
\label{eq:flowmap}
\end{equation}
Thus \(\Phi_\lambda\) is the substrate transport generated by \(U^A\), while \(S\) provides a monotone ordering of that transport.

\begin{remark}
At this stage \(U^A\) is a substrate variable, not yet an independently observed low-energy preferred-frame field. The active exact-closure line is preserved precisely because the present manuscript defines only the bridge architecture, not a surviving extra low-energy sector.
\end{remark}

\section{Local Experiential Space and S-Time}

Given the unit flow field \(U^A\), define the projector orthogonal to the ordered motion by
\begin{equation}
P_{AB} = G_{AB} - U_A U_B,
\label{eq:substrateprojector}
\end{equation}
where \(U_A = G_{AB} U^B\).

\begin{proposition}
The tensor \(P_{AB}\) is positive semidefinite on \(T\Sub\), has kernel spanned by \(U^A\), and is positive definite on the three-dimensional subspace orthogonal to \(U^A\).
\end{proposition}

\begin{proof}
For any tangent vector \(V^A\), decompose
\[
V^A = \alpha U^A + W^A,
\qquad
U_A W^A = 0.
\]
Then
\[
P_{AB} V^A V^B
=
G_{AB} W^A W^B \ge 0,
\]
with equality if and only if \(W^A = 0\), that is, if \(V^A\) is proportional to \(U^A\). Since \(G_{AB}\) is positive definite, its restriction to the orthogonal complement of \(U^A\) is positive definite. The orthogonal complement is three-dimensional because \(\dim \Sub = 4\). Therefore \(P_{AB}\) has the stated properties.
\end{proof}

Equation \eqref{eq:substrateprojector} is the local substrate realization of observed three-dimensional space as a local experiential slice. The word \emph{local} matters. The projector \(P_{AB}\) always defines a local three-space orthogonal to the ordered motion, but those local spaces need not integrate into a single global foliation.

To see the distinction, define the projected antisymmetric derivative
\begin{equation}
\Omega_{AB}
\equiv
P_A{}^C P_B{}^D \nabla_{[C} U_{D]},
\label{eq:substratevorticity}
\end{equation}
where \(\nabla_A\) is the Levi-Civita connection of \(G_{AB}\).

\begin{proposition}
The local experiential three-spaces orthogonal to \(U^A\) integrate to hypersurfaces if and only if
\begin{equation}
U_{[A}\nabla_B U_{C]} = 0,
\label{eq:substratefrobenius}
\end{equation}
equivalently if and only if \(\Omega_{AB}=0\).
\end{proposition}

\begin{proof}
This is the Frobenius criterion applied to the one-form \(U_A\). Hypersurface orthogonality is equivalent to \(U_{[A}\nabla_B U_{C]}=0\). Because \(\Omega_{AB}\) is the orthogonally projected antisymmetric derivative, it vanishes exactly when the orthogonal distribution defined by \(P_{AB}\) is integrable.
\end{proof}

This result is structurally important. It shows that the phrase \emph{local experiential slice} should not be read as a promise of a universal global slicing in every regime. In particular, the derivational bridge remains compatible with the later rotational and frame-dragging sectors, which require nonzero vorticity in appropriate situations.

The order scalar \(S\) then provides the minimal explicit anchor for S-time. A simple kinematic class of S-time functionals is
\begin{equation}
\Order[\Gamma,\Psi]
=
\int_{\Gamma} \rho(X,\Psi)\, dS,
\qquad
\rho(X,\Psi) > 0,
\label{eq:stimefunctional}
\end{equation}
where \(\Gamma\) is an admissible observer history and \(\Psi\) denotes the physical state registered along that history.

For a parameterized admissible history \(X^A(\lambda)\), \eqref{eq:stimefunctional} gives
\begin{equation}
\dv{\Order}{\lambda}
=
\rho\bigl(X(\lambda),\Psi(\lambda)\bigr)\,
\dv{S}{\lambda}.
\label{eq:stimonotone}
\end{equation}
Whenever \(dS/d\lambda > 0\), the S-time functional is monotone. This is the explicit bridge between the substrate order variable and the experienced order of change.

\section{Kinematic Bridge to a Lorentzian Metric}

The next step is the central bridge ansatz of the paper.

\begin{definition}[Bridge metric]
Given the substrate metric \(G_{AB}\) and the unit ordered-flow field \(U^A\), define the bridge metric
\begin{equation}
\BridgeMetric_{AB}
=
P_{AB} - U_A U_B
=
G_{AB} - 2 U_A U_B.
\label{eq:bridgemetric}
\end{equation}
\end{definition}

\begin{proposition}
The bridge metric \(\BridgeMetric_{AB}\) has Lorentzian signature \((-,+,+,+)\).
\end{proposition}

\begin{proof}
For any \(V^A = \alpha U^A + W^A\) with \(U_AW^A=0\),
\begin{equation}
\BridgeMetric_{AB} V^A V^B
=
-\alpha^2 + P_{AB} W^A W^B.
\label{eq:bridgequadratic}
\end{equation}
By the previous proposition, \(P_{AB}\) is positive definite on the three-dimensional space orthogonal to \(U^A\). Therefore \eqref{eq:bridgequadratic} has one negative direction, carried by \(U^A\), and three positive directions, carried by the orthogonal complement. Hence \(\BridgeMetric_{AB}\) is Lorentzian with signature \((-,+,+,+)\).
\end{proof}

Equation \eqref{eq:bridgequadratic} already gives the causal classification induced by the bridge:
\begin{enumerate}
    \item \(V^A\) is timelike if \(\alpha^2 > P_{AB}W^A W^B\);
    \item \(V^A\) is null if \(\alpha^2 = P_{AB}W^A W^B\);
    \item \(V^A\) is spacelike if \(\alpha^2 < P_{AB}W^A W^B\).
\end{enumerate}
This is the first precise statement of how a Lorentzian observer-accessible sector can arise from explicit substrate kinematics while the underlying substrate metric remains positive definite.

Define the corresponding effective four-velocity field on the substrate by
\begin{equation}
\mathfrak u^A = c\, U^A.
\label{eq:substratefourvelocity}
\end{equation}
Then
\begin{equation}
\BridgeMetric_{AB}\,\mathfrak u^A \mathfrak u^B = -c^2,
\label{eq:bridgenorm}
\end{equation}
and the \(\BridgeMetric\)-rest-space projector is
\begin{equation}
\BridgeMetric_{AB}
+
\frac{1}{c^2}\mathfrak u_A \mathfrak u_B
=
P_{AB}.
\label{eq:bridgeprojector}
\end{equation}

Equation \eqref{eq:bridgeprojector} is the key bridge identity. It shows that the local experiential three-space determined on the substrate side is exactly the same object as the observer rest-space geometry defined after the Lorentzian bridge is introduced.

\section{Observer Map and the Effective Relativistic Sector}

To connect the bridge metric to the effective relativistic variables used in the rest of the manuscript sequence, introduce a local observer map
\[
\Xi:M \to \Sub,
\]
where \(M\) is the observer-accessible spacetime manifold and \(\Xi\) is taken locally to be a diffeomorphic identification.

The effective spacetime metric is then defined by pullback:
\begin{equation}
g_{\mu\nu}
=
\Xi^* \BridgeMetric_{AB}
=
\pdv{X^A}{x^\mu}\pdv{X^B}{x^\nu}\BridgeMetric_{AB}\bigl(X(x)\bigr).
\label{eq:effectivepullback}
\end{equation}
Likewise the effective observer field is
\begin{equation}
u^\mu = (\Xi^{-1})_* \mathfrak u^A.
\label{eq:observerpullback}
\end{equation}

\begin{proposition}
The local rest-space metric seen by observers in the effective relativistic sector is the pullback of the substrate projector:
\begin{equation}
h_{\mu\nu}
\equiv
g_{\mu\nu}
+
\frac{1}{c^2}u_\mu u_\nu
=
\Xi^* P_{AB}.
\label{eq:restspacepullback}
\end{equation}
\end{proposition}

\begin{proof}
By \eqref{eq:effectivepullback} and \eqref{eq:observerpullback}, the pair \((g_{\mu\nu},u^\mu)\) is the image of \((\BridgeMetric_{AB},\mathfrak u^A)\) under the same local identification \(\Xi\). Pulling back \eqref{eq:bridgeprojector} therefore gives
\[
\Xi^*\!\left(
\BridgeMetric_{AB}
+
\frac{1}{c^2}\mathfrak u_A \mathfrak u_B
\right)
=
\Xi^* P_{AB},
\]
which is exactly \eqref{eq:restspacepullback}.
\end{proof}

This proposition is the cleanest statement of the derivational bridge. Observed three-dimensional space as a local experiential slice is not left as metaphor. It is identified explicitly with the pullback of the substrate rest-space projector orthogonal to the ordered flow.

The standard operational relativistic relations then follow on the effective side:
\begin{equation}
d\tau^2
=
-\frac{1}{c^2} g_{\mu\nu}\, dx^\mu dx^\nu
=
-\frac{1}{c^2} \BridgeMetric_{AB}\, dX^A dX^B,
\label{eq:bridgepropertime}
\end{equation}
and for null propagation,
\begin{equation}
0
=
g_{\mu\nu}\, dx^\mu dx^\nu
=
\BridgeMetric_{AB}\, dX^A dX^B.
\label{eq:bridgenull}
\end{equation}

\begin{remark}
The congruence-based flow geometry developed elsewhere in the active sequence now acquires a substrate-level origin. After the bridge map is made, the effective observer field \(u^\mu\) is the relativistic image of the deeper ordered-flow field \(U^A\). Expansion, acceleration, vorticity, and shear in the effective geometry are therefore not inserted by hand; they are the observer-accessible kinematic face of the deeper substrate motion.
\end{remark}

\section{What a Completed Derivation Must Achieve}

The present manuscript is still only kinematic unless it is paired with a substrate action. The general derivational target therefore begins with an action of the form
\begin{equation}
S_{\Sub}
=
\int_{\Sub} d^4X\,\sqrt{G}\,
\mathcal{L}_{\Sub}\!\left(
G_{AB},
U^A,
S,
Q,
\psi,
\nabla U,
\nabla S,
\nabla Q,
\ldots
\right),
\label{eq:substrateactionbridge}
\end{equation}
subject to the unit constraint \eqref{eq:unitflow} and monotonicity condition \eqref{eq:ordermonotone}. Here \(\psi\) denotes the matter sector that must eventually be recovered in observer-accessible form.

After solving the substrate equations, eliminating constrained variables, and passing to the long-wavelength observer-accessible description, the derivational bridge would have to yield an effective action of the form
\begin{equation}
S_{\mathrm{eff}}[g,\psi]
=
\frac{c^3}{16\pi G_N}\int d^4x\,\sqrt{-g}\,R
+
S_{\mathrm{matter}}[g,\psi]
+
\Delta S[g,\psi].
\label{eq:bridgeeffectiveaction}
\end{equation}
In the active exact-closure line, successful recovery requires
\begin{equation}
\Delta S[g,\psi] = 0
\label{eq:deltazero}
\end{equation}
at the effective level. If \(\Delta S\neq 0\) survives as an observable low-energy correction, the result is no longer the exact-closure theory already established elsewhere in the sequence.

The resulting derivational success conditions can now be stated precisely.
\begin{enumerate}
    \item \textbf{Bridge closure.} The map \((G_{AB},U^A,S,Q)\mapsto g_{\mu\nu}\) defined through \eqref{eq:bridgemetric} and \eqref{eq:effectivepullback} must be regular, unique in the observer-accessible regime, and free of ambiguities that would introduce multiple competing low-energy metrics.
    \item \textbf{Universal matter coupling.} The observer-accessible matter sector must couple universally to the same \(g_{\mu\nu}\) that governs gravity, null propagation, and proper time.
    \item \textbf{Infrared Einsteinian recovery.} The long-wavelength metric dynamics must reduce exactly to the Einstein equations with ordinary matter coupling, equivalently to \eqref{eq:bridgeeffectiveaction} with \(\Delta S=0\).
    \item \textbf{Mode control.} Any extra substrate modes carried by \(U^A\), \(S\), or \(Q\) must either be constrained, screened, or decoupled so that no independent low-energy preferred-frame observable sector survives inside \TheoryName{}.
    \item \textbf{Compatibility with the active sequence.} The recovered observer field \(u^\mu\) and rest-space geometry \(h_{\mu\nu}\) must match the congruence-based flow geometry, weak-field structure, causal structure, and local SR recovery already fixed in the exact-closure manuscripts.
\end{enumerate}

These criteria are intentionally strong. They ensure that the derivational program is answerable to the exact-closure benchmark rather than drifting into a looser metaphorical or deformation-first narrative.

\section{What This Manuscript Establishes and What It Does Not}

The present manuscript establishes the following positive structure.
\begin{enumerate}
    \item It makes the substrate variables explicit through \((\Sub,G_{AB},U^A,S,Q)\).
    \item It identifies the local experiential three-space with the substrate projector \(P_{AB}\) orthogonal to the ordered flow.
    \item It gives a concrete Lorentzian bridge metric \(\BridgeMetric_{AB}=G_{AB}-2U_AU_B\).
    \item It defines the observer map by which the effective metric \(g_{\mu\nu}\), the effective observer field \(u^\mu\), and the rest-space metric \(h_{\mu\nu}\) are recovered.
    \item It states exact derivational success conditions in a form that can be tested by any later substrate action.
\end{enumerate}

The present manuscript does \emph{not} establish the following.
\begin{enumerate}
    \item It does not derive the Einstein equations from a completed substrate action.
    \item It does not prove that a specific choice of \(\mathcal{L}_{\Sub}\) reproduces \eqref{eq:bridgeeffectiveaction} with \(\Delta S=0\).
    \item It does not complete the mode counting or stability analysis of a concrete substrate model.
    \item It does not introduce any established non-GR low-energy observable signature within \TheoryName{}.
\end{enumerate}

\begin{remark}
The correct scientific status of the manuscript is therefore intermediate but substantive. It closes the missing kinematic bridge from ontology to an explicit candidate relativistic sector, while leaving the decisive dynamical proof to later work. Until that later work succeeds, exact closure remains the benchmark and reversion theory of the framework.
\end{remark}

\section{Conclusion}

The project could not advance indefinitely by repeating exact GR / SR recovery statements. What it needed next was a disciplined bridge from ontology to explicit substrate kinematics. The present manuscript supplies that bridge.

The core move is simple but consequential. The \Aether{} is represented by a four-dimensional substrate manifold with positive-definite metric \(G_{AB}\). The \Aflow{} is represented kinematically by a unit ordered-flow field \(U^A\) together with an order scalar \(S\) monotone along that flow. Observed three-dimensional space is represented by the local projector orthogonal to \(U^A\). The observer-accessible Lorentzian sector is then obtained through the bridge metric \(\BridgeMetric_{AB}=G_{AB}-2U_AU_B\) and its pullback to the effective spacetime metric \(g_{\mu\nu}\).

This is not yet a completed derivation of Einsteinian gravity. It is the missing derivational architecture on which such a derivation would have to stand. The exact-closure theory of \TheoryName{} therefore remains scientifically primary, but the path beyond exact closure is now sharper: any later substrate action must answer to the bridge stated here and recover the exact relativistic benchmark already fixed by the active sequence.

\end{document}
